\subsection{مبانی نظری و فرمول‌بندی ریاضی}
فرض کنید ماتریس وزن هدف در لایه‌های تصویرسازی توجه یا پروجکشن‌های خطی به صورت $\mathbf{W} \in \mathbb{R}^{d_{\text{out}} \times d_{\text{in}}}$ داده شده باشد. الگوریتم \lr{Delta-KV} ماتریس وزن بازسازی‌شده را به فرم مجموع رتبه پایین و مانده تنک مدل‌سازی می‌کند:
\begin{equation}
\mathbf{W}_{\text{comp}} = \mathbf{L} + \mathbf{R}
\end{equation}
به طوری که:
\begin{enumerate}
    \item مؤلفه با رتبه پایین ساختاری $\mathbf{L}$ با استفاده از تجزیه مقادیر منفرد بر روی ماتریس وزن به دست می‌آید:
    \begin{equation}
    \mathbf{W} = \mathbf{U} \mathbf{\Sigma} \mathbf{V}^\top, \quad \mathbf{L} = \mathbf{U}_{:, :r} \mathbf{\Sigma}_{:r} \mathbf{V}_{:, :r}^\top
    \end{equation}
    که در آن $r = \max(1, \min(r_{\text{target}}, d_{\text{in}}, d_{\text{out}}))$ رتبه فشرده‌سازی انتخاب‌شده است.
    \item ماتریس مانده خام برابراست با تفاضل ماتریس اصلی و تقریب کم‌مرتبه:
    \begin{equation}
    \mathbf{E} = \mathbf{W} - \mathbf{L}
    \end{equation}
    سپس بر مبنای نرخ تنک‌پذیری تعیین‌شده $s \in [0, 1]$، عملگر گزینش بزرگ‌ترین قدرمطلق‌ها (\lr{Top-k Magnitude Operator}) بر روی المان‌های برداری‌شده $|\mathbf{E}|$ اعمال می‌شود. با تعیین آستانه برش $\tau = \text{Quantile}(|\mathbf{E}|, 1 - s)$، ماتریس مانده تنک $\mathbf{R}$ محاسبه می‌گردد:
    \begin{equation}
    \mathbf{R}_{i, j} = \begin{cases} 
    \mathbf{E}_{i, j} & \text{if } |\mathbf{E}_{i, j}| \ge \tau \\
    0 & \text{otherwise}
    \end{cases}
    \end{equation}
\end{enumerate}

خطای بازسازی فروبنیوس به صورت زیر اعتبارسنجی می‌شود:
\begin{equation}
\epsilon_{\text{rel}} = \frac{\|\mathbf{W} - (\mathbf{L} + \mathbf{R})\|_F}{\max(\|\mathbf{W}\|_F, 10^{-12})}
\end{equation}
در صورتی که $\epsilon_{\text{rel}} > \epsilon_{\text{max}}$ باشد، الگوریتم لایه را به ساختار متراکم اولیه بازمی‌گرداند. علاوه بر این، در حالت اجرای برخط کش، ماژول شامل کلاس ذخیره‌ساز پویا \lr{DeltaKVStore} است که با سنجش شباهت کسینوسی بردار پنهان جاری با حافظه تاریخی $K/V$ مراجع، تفاضل بردار را در قالب دلتای فشرده بازنمایی می‌نماید:
\begin{equation}
\mathbf{k}_{\text{residual}} = \mathbf{k} - \mathbf{r}_k, \quad \hat{\mathbf{k}} = \mathbf{r}_k + \mathbf{k}_{\text{residual}} \odot \mathbb{I}(\|\mathbf{k}_{\text{residual}}\|_2 > \theta_{\text{res}})
\end{equation}
